\section{Common routing and removal of invariant occupation}\label{sec:occupation-proof}
We prove Theorem~\ref{thm:main}(i) without using any realization or any metric geometry.

\begin{lemma}[A joint routing separation criterion]\label{lem:routing}
Fix common $p$ and vectors $z_s^\theta\in C$. There exist common Borel routing probabilities \eqref{eq:routing} such that, for all $\theta$,
\begin{align}
 z_i^\theta-a_i^\theta
  &=\sum_{s,a}p_{sa}\int\gamma_i(s,a,y)L_a^\theta(y)z_s^\theta\,\lambda_a(dy)
       &&(i\in S),\label{eq:balance-live}\\
 d_j^\theta
  &=\sum_{s,a}p_{sa}\int\gamma_j(s,a,y)L_a^\theta(y)z_s^\theta\,\lambda_a(dy)
       &&(j\in J)\label{eq:balance-terminal}
\end{align}
if and only if \eqref{eq:joint-support} holds for every finite $F$ and every displayed dual vector.
\end{lemma}
\begin{proof}
For each pair $(s,a)$ the simplex-valued kernels form a weak-star compact convex subset of $L^\infty(\lambda_a)^{S\sqcup J}$. Positivity and the sum-one constraint are weak-star closed: test them against nonnegative $L^1$ functions. The product over finitely many $(s,a)$ is compact and convex. Integration against $L_a^\theta(y)z_s^\theta$ is weak-star continuous by \eqref{eq:raw-instrument}. For a finite $F$, the vector of the right sides of \eqref{eq:balance-live}--\eqref{eq:balance-terminal} therefore has a compact convex image.

Its support function is the right side of \eqref{eq:joint-support}. Indeed, after summing the paired contributions over $\theta\in F$, one maximizes a linear function over a probability simplex, for each $(s,a,y)$. A maximizing label can be selected by the smallest-index rule among finitely many measurable functions. This gives a Borel maximizer and the stated integral of a maximum. Finite-dimensional separation proves that the target vector belongs to the image exactly when all support inequalities hold.

For an arbitrary $\Theta$, impose the balance equations first on each finite $F$. Their solution sets in the same weak-star compact routing space are closed and have the finite intersection property: the union of finitely many finite model sets is finite. A common routing kernel exists by compactness. On standard Borel report spaces its equivalence classes admit Borel versions. Changing representatives on a common $\lambda_a$-null set and assigning a fixed destination there makes \eqref{eq:routing} pointwise valid without changing any model law. Necessity is integration of the pointwise simplex bound.
\end{proof}

For such a kernel define positive linear maps on ensembles,
\[
 (P_\theta z)_i=\sum_{s,a}p_{sa}\int\gamma_i(s,a,y)L_a^\theta(y)z_s\,\lambda_a(dy),
 \quad
 (Q_\theta z)_j=\sum_{s,a}p_{sa}\int\gamma_j(s,a,y)L_a^\theta(y)z_s\,\lambda_a(dy).
\]
Write $\boldsymbol e(z)=\sum_s e(z_s)$, with the analogous convention on terminal ensembles. Normalization gives
\begin{equation}\label{eq:mass-conservation}
 \boldsymbol e(P_\theta z)+\boldsymbol e(Q_\theta z)=\boldsymbol e(z)
 \quad\text{for every positive live ensemble }z.
\end{equation}

\begin{lemma}[Budget-preserving removal]\label{lem:removal}
Let $P,Q$ be positive ensemble maps satisfying \eqref{eq:mass-conservation}. If a finite positive ensemble solves $z=a+Pz$, then
\begin{equation}\label{eq:removal}
 h=\lim_{n\to\infty}P^nz\in C^S,
 \qquad Ph=h,\qquad Qh=0,
 \qquad z^{\mathrm{act}}=\sum_{n=0}^\infty P^na=z-h.
\end{equation}
The common observer started at preparation $a$ has actual continuing ensemble $z^{\mathrm{act}}$, terminal ensemble $Qz$, and mean acquisition $\boldsymbol e(z^{\mathrm{act}})\le\boldsymbol e(z)$. Every positive linear running-cost functional is no larger on the actual ensemble.
\end{lemma}
\begin{proof}
Since $Pz=z-a\le_C z$, the sequence $P^nz$ decreases in the cone order. Its scalar mass decreases to a finite limit. Compactness of the base implies an estimate $\norm w\le C_D\boldsymbol e(w)$ for positive ensembles. Thus, for $n\ge m$,
\[
 \norm{P^mz-P^nz}\le C_D\bigl(\boldsymbol e(P^mz)-\boldsymbol e(P^nz)\bigr),
\]
so the sequence converges in the finite-dimensional space. Its limit $h$ is positive, and continuity of $P$ gives $Ph=h$. Applying \eqref{eq:mass-conservation} to $h$ gives $\boldsymbol e(Qh)=0$. Strict positivity of the mass functional on the cone implies $Qh=0$, component by component.

Iteration of $z=a+Pz$ yields
\[
 z=\sum_{n=0}^{t-1}P^na+P^tz.
\]
Taking the limit proves \eqref{eq:removal}. Starting the physical experiment at $a$, induction with the raw instruments shows that $P^na$ is exactly the unnormalized predictive ensemble of surviving live labels after $n$ acquisitions. This statement uses only linearity before normalization and the fact that the same report-routing kernel is executed at every step. Tonelli's theorem, applied first to masses and then to the bounded base coordinates, gives the actual occupation and $\E\tau=\boldsymbol e(z-h)<\infty$. Hence $\tau<\infty$ almost surely. Its terminal ensemble is $Q\sum_nP^na=Q(z-h)=Qz$. A positive linear cost decreases when the positive component $h$ is removed.
\end{proof}

\begin{example}[Why the candidate is not the acquired law]\label{ex:phantom}
Take the scalar cone $\R_+$, live states $0,1$, terminal state $2$, preparation $a=(1,0)$. State $0$ always halts after one call; state $1$ always returns to itself. Then
\[
 P=\begin{pmatrix}0&0\\0&1\end{pmatrix},\qquad Q=(1,0),\qquad z=(1,h),\quad d=1
\]
solves the balance equations for every $h\ge0$. The prepared run never visits state $1$. Its mean acquisition is $1$, not $1+h$. The extra component changes neither the terminal law nor an upper budget constraint. This is the elementary phantom-measure phenomenon in occupation theory~\cite{DPR}; its explicit role here is to delimit the interpretation of a predictive certificate.
\end{example}

\begin{proof}[Proof of Theorem~\ref{thm:main}(i)]
For an admissible observer, linear conditional integration and Tonelli give \eqref{eq:balance-live}--\eqref{eq:balance-terminal} with the actual ensembles \eqref{eq:actual-occupation}. The mean budget gives \eqref{eq:budget-candidate}; Lemma~\ref{lem:routing} gives \eqref{eq:joint-support}. Conditional expectation at the protected audit gives risk $\sum_jG_\theta(d_j^\theta,r_j)$. Therefore $\Phi^\Theta\le\calR_\aut^\Theta$.

Conversely, a feasible candidate has one common routing witness by Lemma~\ref{lem:routing}. For each $\theta$, Lemma~\ref{lem:removal} identifies the actual ensemble generated by its preparation. The same kernel, not a model-specific choice of kernel, is used for every $\theta$. Its actual expected acquisition is at most $N$, and its terminal ensemble is exactly $d^\theta$. Its loss is therefore the candidate loss. This proves the reverse inequality, including at the level of infima.

It remains to show attainment. For fixed cardinalities, action rows lie in compact simplices and decoders in $H^J$. For each model, the positive ensemble sets in \eqref{eq:budget-candidate} are compact, because their masses are bounded and $D$ is compact. Take the product of these compact sets over $\Theta$. Each inequality \eqref{eq:joint-support} depends on finitely many models and is continuous in $p,z,d$. To verify continuity of its integral, bound the difference of maxima by the largest difference of the finitely many linear terms, and use the integrability in \eqref{eq:raw-instrument}. The feasible set is therefore compact. It is nonempty by the one-call stopping program. The objective, a supremum of continuous model losses, is lower semicontinuous; it attains its minimum. There are only finitely many cardinality choices. The implementing kernel supplied above attains the resulting global minimum.

For \eqref{eq:finite-models}, the inequality from right to left follows by restriction. Suppose that $r$ is strictly smaller than the full-family optimum and no finite subfamily has optimum greater than $r$. Work in the finite disjoint union of the compact candidate spaces, adding the closed constraints $\sum_jG_\theta(d_j^\theta,r_j)\le r$. Any finite collection of support inequalities and risk constraints refers to finitely many models. A minimizing candidate for their union satisfies the collection. Coordinates for other models may be filled arbitrarily in their nonempty budget sets. Thus these closed sets have the finite intersection property and yield a full-family feasible candidate of risk at most $r$, a contradiction. This also proves the strict finite-model obstruction. There is no interchange of $\inf$ and $\sup$, and no bound on the size of the obstructing finite set is asserted.
\end{proof}

\begin{corollary}[Additional nonnegative costs]\label{cor:costs}
In Theorem~\ref{thm:main}(i), finitely many further constraints of the form
\[
 \sum_{s,a}p_{sa}\,\kappa_{\ell,sa}^\theta(z_s^\theta)\le B_\ell,
\]
where each $\kappa_{\ell,sa}^\theta$ is a fixed nonnegative linear functional representing an actual conditional expected primitive cost, can be added. The unit acquisition bound is retained. The value equality and finite-model statement remain valid, and the value is attained whenever the constrained admissible class is nonempty.
\end{corollary}
\begin{proof}
The additional inequalities are closed. After extracting a common routing witness, removal of its invariant occupation can only decrease each displayed cost. The compactness and finite-intersection arguments are unchanged.
\end{proof}

\begin{remark}[Necessity versus a finite computation]
The support inequalities are an exact realizability criterion, not a finite numerical test in general. Infeasibility for a fixed candidate and finite model family has a strict finite-dimensional separating vector. The global finite-model conclusion is stronger but still non-effective: it supplies neither a uniformly bounded witness size nor an algorithm for locating it. These distinctions prevent a compactness proof from being reported as a finite-bit synthesis theorem.
\end{remark}
